\section{Related Work}
\label{sec:related_work}

Our work is positioned at the intersection of static model merging, test-time adaptation, and dynamic parameter routing. In this section, we review each area and highlight the methodological and paradigm differences that our work addresses.

\subsection{Static Model Merging}
Parameter fusion or model merging has emerged as a key approach to consolidate multiple specialized networks fine-tuned from a single pre-trained initialization \cite{task_arithmetic, ties_merging, dare}. The foundational protocol, Task Arithmetic \cite{task_arithmetic}, demonstrates that task-specific weights can be combined via a simple linear combination of their task vectors: $V_k = W_k - W_{base}$. However, as the number of tasks grows, linear interpolation suffers from severe interference due to parameter conflicts. 
To address this, TIES-Merging \cite{ties_merging} introduces sign consensus and magnitude thresholding to prune redundant, low-magnitude parameter changes before merging. Similarly, DARE \cite{dare} applies random delta-pruning and scaling to maintain output distribution consistency. Recently, Sparse Task Arithmetic (STA) \cite{sparse_task_arithmetic} challenges the sign-consensus heuristics of TIES-Merging and DARE, arguing that sign-conflict resolution is redundant once update under-scaling is corrected.
Despite these advances, static merging models are fundamentally limited: they compress task vectors into a single set of static, input-agnostic weights, which inevitably results in trade-offs and performance degradation (representation collapse) on high-conflict datasets.

\subsection{Test-Time Adaptation vs.\ Offline Calibration}
To overcome the limitations of static weights without the overhead of joint multi-task training, researchers have explored dynamic model-merging adaptations. A prominent paradigm is **Test-Time Adaptation (TTA)**, which dynamically tunes merging coefficients at inference time. For example, AdaMerging \cite{adamerging} optimizes merging coefficients on-the-fly by running active gradient descent to minimize the entropy of model predictions on incoming, unlabeled test streams. 
While elegant, online TTA is computationally expensive. It requires backpropagation passes during inference, introduces substantial latency, and is highly susceptible to stream noise, temporal order, and local label shift (the "Overfitting-Optimizer Paradox"). Under highly non-i.i.d.\ streams, unsupervised test-time optimization can cause coefficients to drift catastrophically, degrading performance.

In contrast, **Offline-Calibrated Dynamic Routing** (such as QWS-Merge or our proposed routers) optimizes its routing heads or scaling parameters prior to deployment on a tiny, balanced calibration set. At inference time, these routing heads execute as a pure, lightweight forward pass with zero active optimization, zero backward passes, and zero latency overhead. 
While prior work frequently compares TTA and offline calibration on a simple accuracy basis, we argue from a methodological perspective that they represent fundamentally different computational and deployment paradigms. Offline calibration is operationally far more practical for real-time, low-latency deployment. In this work, we adopt the offline, supervised calibration protocol to guarantee a fair, robust, and reproducible evaluation of all trainable routing methods.

\subsection{Dynamic Parameter Routing and Complex Metaphors}
The limitations of both static merging and online test-time adaptation have led to the development of dynamic model merging, where parameters are dynamically assembled for each input batch or sample via routing heads. A standard baseline in this space is the classical Linear Router, which projects pooled intermediate representations to task coefficients using a simple linear projection followed by Softmax.

However, recent publications have shown a troubling trend of introducing highly complex, over-engineered mathematical metaphors to replace this classical baseline. The most prominent example is Quantum Wavefunction Superposition Merging (QWS-Merge) \cite{qws_merge}, which models task-specific parameters as quantum eigenstates in a Hilbert space. QWS-Merge maps input representations to a unit sphere and uses cosine wave-interference equations to compute routing weights. The authors claim that this "quantum phase projection" is essential to avoid the representation collapse that causes classical linear routers to fail on high-conflict tasks like SVHN.

From a methodological standpoint, we argue that the deep learning literature frequently suffers from a lack of rigorous baseline optimization. Often, a new, complex method is compared against an under-tuned or un-regularized version of a classical baseline, leading to false progress. By applying Occam's razor, we demonstrate that when the classical Linear Router is regularized to prevent over-scaling, or when it is allowed to utilize layer-wise scaling parameters, it completely deconstructs the necessity of the quantum wavefunction metaphor.
